\chapter{Compresia imaginilor: JPEG versus JPEG2000}
\label{ch:compresie}

Comprimarea unei imagini se reduce, în ambele standarde, la aceeași idee de
la finalul capitolului trecut: transformă, păstrează coeficienții mari, taie
sau cuantizează grosier restul. JPEG transformă cu DCT pe blocuri mici.
JPEG2000 transformă cu DWT, pe toată imaginea, folosind exact algoritmul lui
Mallat.
Acest capitol arată pipeline-ul JPEG pas cu pas, apoi pune cele două
transformate față în față, la bitrate egal, cu cifre măsurate.

\section{Pipeline-ul JPEG}

Baseline JPEG parcurge șase etape: conversie RGB în YCbCr (separă luminanța
de crominanță), subeșantionare de crominanță (ochiul distinge culoarea mai
slab decât luminozitatea), DCT pe blocuri de $8\times 8$, cuantizare cu
tabelul Annex K scalat după regula IJG, apoi zigzag și codare pe lungimi de
serie. Demo-ul din tur (Figura~\ref{fig:jpeg-pipeline}) rulează toate șase
pe o fotografie reală, nu pe un gradient sintetic, un gradient neted n-are
ce pierde la subeșantionare, deci ar ascunde exact etapa pe care vrea s-o
arate.

\begin{figure}[htbp]
\centering
\includegraphics[width=\textwidth]{44-jpeg-pipeline.png}
\caption{Prima etapă a pipeline-ului JPEG pe o fotografie reală: separarea
  Y/Cb/Cr, cu formulele de conversie folosite efectiv de back-end.}
\label{fig:jpeg-pipeline}
\end{figure}

Cele două standarde lasă în urmă tipuri diferite de artefact la compresie
mare. JPEG, cu blocuri independente, lasă blocare: grila de $8\times 8$
devine vizibilă. JPEG2000, cu o transformată globală, lasă
\emph{ringing}~\cite{taubmanmarcellin2002}: oscilații fine în jurul
muchiilor, nu o grilă. JPEG e progresiv doar dacă fișierul e codat explicit
în acel mod din T.81, baseline-ul, cel mai comun pe web, se decodează
secvențial~\cite{itut81}. JPEG2000 capătă scalabilitate progresivă direct
din construcția multi-rezoluție: aceeași descompunere care dă comprimarea dă
și un flux din care poți citi mai întâi o versiune la rezoluție mică, apoi
detaliile~\cite{taubmanmarcellin2002}.

\section{O comparație la bitrate egal}
\label{sec:comparatie}

O primă versiune a acestei comparații compara cele două codecuri la un
parametru de „calitate” egal, calibrat diferit pe fiecare parte, partea DCT
folosea tabelul real Annex K, partea wavelet un prag liniar ad-hoc, fără
nicio legătură cu costul în biți, și scotea un rezultat înșelător, cu JPEG
câștigând la orice calitate testată. Comparația de mai jos egalează
bitrate-ul, nu eticheta de calitate, și măsoară amândouă codecurile cu
același estimator de cost.

\paragraph{Costul în biți.} Fiecare codec e măsurat cu entropia de ordin 0 a
simbolurilor lui cuantizate, înmulțită cu numărul lor,
\[
R = N \cdot H(\mathbf{p}), \qquad H(\mathbf{p}) = -\sum_i p_i \log_2 p_i,
\]
unde $p_i$ e frecvența empirică a simbolului $i$ în histograma
coeficienților cuantizați. Un JPEG sau un JPEG2000 adevărat iese mai mic
decât atât, pentru că folosește o codare entropică reală. Huffman,
aritmetică sau, la JPEG2000, EBCOT cu context, dar cum ambele părți sunt
măsurate cu același estimator, raportul dintre cele două costuri rămâne
comparabil, chiar dacă niciuna dintre cifrele absolute nu e dimensiunea unui
fișier real.

\paragraph{Partea DCT.} Imaginea se taie în blocuri de $8\times 8$ pixeli, iar
fiecare bloc trece prin transformata de cosinus discretă ortonormală,
\[
X_{u,v} = \tfrac14 C_u C_v \sum_{i=0}^{7}\sum_{j=0}^{7} x_{i,j}
  \cos\!\left[\frac{(2i+1)u\pi}{16}\right]
  \cos\!\left[\frac{(2j+1)v\pi}{16}\right], \qquad
C_0 = \tfrac{1}{\sqrt2},\ C_{k>0} = 1,
\]
care scrie blocul ca sumă de $64$ tipare de cosinus, de la cel constant la
cel mai alternant. Urmează cuantizarea cu tabelul de luminanță Annex~K
scalat după regula IJG,
\[
\text{scală} = \begin{cases} 5000/\text{calitate} & \text{calitate} < 50 \\ 200 - 2\cdot\text{calitate} & \text{calitate}\ge 50 \end{cases}, \qquad
Q = \mathrm{clip}\!\left(\left\lfloor \frac{Q_{50}\cdot\text{scală}+50}{100}\right\rfloor,\, 1,\, 255\right).
\]

\paragraph{Partea wavelet primește bugetul de biți al părții DCT.} În loc de
propriul parametru de calitate, partea wavelet caută prin bisecție
geometrică pasul de cuantizare $\Delta$ cel mai fin care încă încape în
bugetul de biți al părții DCT (treizeci de iterații, interval de căutare
$[0.02,\,8192]$, oprire la o precizie relativă de $0.2\%$). Cuantizarea e cu
\emph{zonă moartă}, cu un pas propriu pe fiecare subbandă $b$,
\[
q_b = \mathrm{sign}(c_b)\left\lfloor \frac{|c_b|}{\Delta / g_b}\right\rfloor, \qquad
\hat c_b = \mathrm{sign}(q_b)\left(|q_b|+\tfrac12\right)\frac{\Delta}{g_b},
\]
cu reconstrucția la \emph{mijlocul} intervalului de cuantizare, convenția
standard JPEG2000. Factorul $g_b$ e câștigul de sinteză al subbenzii:
transformata folosită implicit, \txt{bior4.4}, e chiar filtrul 9/7 al
JPEG2000, construit istoric prin schema de \emph{lifting}~\cite{sweldens1998},
iar fiind biortogonală, nu ortonormală, o eroare de un coeficient nu produce
aceeași eroare de reconstrucție în toate subbenzile. Câștigul $g_b$ se
măsoară direct, nu prin formulă: un impuls unitar plasat în centrul
subbenzii, trecut prin transformata inversă singur, iar norma $L_2$ a
rezultatului e $g_b$, exact tehnica JPEG2000 de pași de cuantizare diferiți
pe subbandă~\cite{taubmanmarcellin2002}, aplicată empiric.

\paragraph{Metrici.} PSNR e definiția uzuală,
$\mathrm{PSNR} = 10\log_{10}(255^2/\mathrm{MSE})$. SSIM e implementarea
proprie a formulei lui Wang și colaboratorii,
\[
\mathrm{SSIM}(x,y) = \frac{(2\mu_x\mu_y+C_1)(2\sigma_{xy}+C_2)}{(\mu_x^2+\mu_y^2+C_1)(\sigma_x^2+\sigma_y^2+C_2)}, \qquad
C_1=(0.01L)^2,\ C_2=(0.03L)^2,\ L=255,
\]
pe ferestre $7\times 7$~\cite{wangbovik2004}, verificată numeric față de
\txt{skimage.metrics.structural_similarity}: $0.861644$ (implementare
proprie) față de $0.861132$ (referință) pe imaginea \txt{kodim01} la
calitate 30, diferența venind din tratarea marginilor ferestrei.

\paragraph{Măsurători.} Tabelul~\ref{tab:compare} arată diferența wavelet
minus DCT, la bitrate egal, pentru \txt{bior4.4} pe patru niveluri, pe trei
imagini standard de test.

\begin{table}[htbp]
\centering
\caption{Diferența wavelet minus DCT la bitrate egal (\txt{bior4.4}, 4 niveluri).}
\label{tab:compare}
\begin{tabular}{@{}llrrr@{}}
\toprule
Imagine & Calitate & bpp & $\Delta$PSNR & $\Delta$SSIM \\
\midrule
kodim01 ($768\times 512$) & 10 & 0.512 & $+0.23$~dB & $-0.026$ \\
kodim01 & 20 & 0.865 & $+0.74$~dB & $-0.007$ \\
kodim01 & 30 & 1.132 & $+1.12$~dB & $+0.003$ \\
kodim01 & 50 & 1.522 & $+1.71$~dB & $+0.010$ \\
kodim01 & 75 & 2.159 & $+2.55$~dB & $+0.012$ \\
\txt{peppers_512} & 10 & 0.537 & $-0.33$~dB & $-0.017$ \\
\txt{peppers_512} & 30 & 1.016 & $+0.42$~dB & $-0.011$ \\
\txt{peppers_512} & 50 & 1.319 & $+0.90$~dB & $-0.005$ \\
\txt{baboon_512} & 10 & 0.484 & $+0.12$~dB & $-0.036$ \\
\txt{baboon_512} & 30 & 1.135 & $+0.96$~dB & $+0.006$ \\
\txt{baboon_512} & 50 & 1.567 & $+1.59$~dB & $+0.015$ \\
\bottomrule
\end{tabular}
\end{table}

\emph{Notă}: în ciuda numelui de fișier, \txt{peppers_512}, \txt{baboon_512},
\txt{house_512} și \txt{lake_512} au de fapt $200\times 200$ pixeli.

Demo-ul din tur (Figura~\ref{fig:compare-demo}) rulează exact această
comparație, live, pe imaginea și calitatea alese: la aproximativ 1.01 biți
pe pixel, wavelet câștigă pe PSNR ($+0.93$~dB) și pierde la limită pe SSIM
($-0.001$), cu ambele metrici afișate pe ecran.

\begin{figure}[htbp]
\centering
\includegraphics[width=\textwidth]{47-compare-demo.png}
\caption{Comparație la bitrate egal ($\approx 1.01$~bpp): wavelet câștigă pe
  PSNR și pierde la limită pe SSIM, cu ambele metrici afișate pe ecran.}
\label{fig:compare-demo}
\end{figure}

Tabelul spune ceva simplu: wavelet câștigă pe PSNR aproape peste tot, cu
marja crescând odată cu bitrate-ul. SSIM îl favorizează ușor pe JPEG la
bitrate foarte mic, și e un rezultat real, nu o eroare de măsurare, tabelul
Annex K e ponderat perceptual, construit din studii asupra sensibilității la
contrast a ochiului uman, în timp ce cuantizatorul wavelet folosit aici e
optim pe eroare medie pătratică, nu pe percepție. Restul avantajului pe care
JPEG2000 îl are în literatură la bitrate mic vine din codarea entropică cu
context, EBCOT~\cite{taubmanmarcellin2002}, pe care estimatorul de entropie
de ordin 0 folosit aici nu îl modelează: comparația izolează transformata și
cuantizarea, nu codecul complet.
